% ===================================================================
\section{Discussion}
\label{sec:discussion}

\subsection{The Bridge as Programmable Substrate}

Paper~3 framed the result as \emph{discovery}: the Steersman
discovers rhombic dodecahedral geometry in the bridge matrix. The
present paper reframes this as \emph{programming}: the Steersman
programs whatever geometry it is told, provided the pair
specification corresponds to a coherent spatial structure.

The evidence supports the reframing. The tesseract result
(Section~\ref{sec:tesseract}) demonstrates that the Steersman
programs 4D geometry with the same fidelity as 3D: four independent
$2 \times 2$ blocks, clean eigenvalue split, coupling ratio
exceeding 41{,}000:1. The octahedral result (O-001) demonstrates
programming at minimum dimensionality, achieving the programme's
strongest block-diagonal signal (473{,}622:1 per-bridge mean). The
spectral attractor (Section~\ref{sec:attractor}) shows that the
bridge has a natural unprogrammed state---connected, isotropic,
Fiedler~$\approx 0.09$---and that contrastive loss adds direction
to this connectivity.

The four-regime taxonomy (Section~\ref{sec:programmability}) sharpens
the framing further. The bridge is a programmable substrate, but the
programming language is \emph{spatial}. The Steersman compiles
polytope-derived pair specifications into block-diagonal structure.
It rejects non-spatial specifications---random partitions and
number-theoretic relationships---by driving the bridge into
connectivity collapse. The selectivity is a feature: it confirms
that successful topology programming encodes genuine geometric
information, not an artifact of the loss formulation.

\subsection{The Selective Geometric Compiler}

The central thesis upgrade from Paper~3 is precise. The Steersman
is not a universal topology enforcer---it does not program arbitrary
pair specifications into block structure. It is a \emph{selective
geometric compiler}: it accepts pair specifications derived from
the co-axial face or vertex structure of polytopes and rejects
all others.

This selectivity is demonstrated by the four-way contrast at $n = 6$:

\begin{itemize}
  \item \textbf{RD pairs (H-ch6):} 70{,}404:1. Spatial
    specification. Block-diagonal.
  \item \textbf{Emanation (E-001):} 1.12:1. Same spatial pairs,
    indirect application through master bridge. Weak coherence but
    no block-diagonal.
  \item \textbf{Wrong-labels (WL-001):} $8.7 \times 10^{-6}$:1.
    Random pairs. Connectivity collapse.
  \item \textbf{Resonance (R-001):} $8.9 \times 10^{-6}$:1.
    Number-theoretic pairs. Connectivity collapse.
\end{itemize}

The selectivity has two dimensions. The first is geometric: the pair
specification must encode a coherent spatial decomposition (Regimes~1
vs.~4). The second is architectural: the contrastive signal must reach
individual bridge matrices directly, not through hierarchical
indirection (Regimes~1 vs.~2).

\subsection{Number-Theoretic Structure Is Insufficient}

The functional equivalence of WL-001 and R-001 (maximum divergence
${<}\,10\%$ across 10{,}000 steps) has a clear implication:
\emph{mathematical structure without spatial embedding carries no
information that the contrastive loss can use}. The Sophie Germain
relationship ($2(11) + 1 = 23$), the twin prime pair $(29, 31)$, and
the residue class partition are genuine algebraic structures. They
produce the same outcome as random noise.

This result constrains the class of viable topology specifications.
The pair specification must encode a property of the bridge's
gradient landscape, not merely a mathematical relationship among
channel indices. Polytope co-axial structure encodes such a
property: it identifies axis-aligned subspaces of the bridge matrix
that can be independently optimized. Number-theoretic relationships
do not identify such subspaces.

\subsection{Spectral-Contrastive Decomposition}

The decomposition of the Steersman into spectral (connectivity) and
contrastive (topology) components has a precise analogue in graph
theory. The Fiedler eigenvalue measures algebraic
connectivity~\cite{fiedler1973}---how well-connected the graph is.
The contrastive loss specifies edge weights---which connections are
strengthened. Spectral regularization without contrastive
specification produces an Erd\H{o}s--R\'enyi-like coupling
pattern (all edges approximately equal); adding contrastive
specification produces a stochastic block model~\cite{holland1983}
with known community structure.

The four regimes map onto this decomposition:

\begin{itemize}
  \item \textbf{Spectral attractor:} Spectral loss alone. Connectivity
    without topology. The bridge graph is an approximately regular graph.
  \item \textbf{Block-diagonal:} Spectral + coherent contrastive. The
    bridge graph is a disjoint union of cliques (the blocks).
  \item \textbf{Hierarchical coherence:} Spectral + indirectly applied
    contrastive. The bridge graph is weakly modular---slight clustering
    without full community separation.
  \item \textbf{Collapse:} Spectral + incoherent contrastive.
    Destructive interference between the two losses drives the graph
    toward the empty graph.
\end{itemize}

This decomposition suggests that the two loss components are
independently tunable but interact non-linearly. When they are
aligned (coherent contrastive + spectral), they produce strong
structure. When they conflict (incoherent contrastive + spectral),
the contrastive loss dominates and suppresses connectivity below the
spectral target. Understanding this interaction is a prerequisite
for designing more sophisticated topology specifications.

\subsection{The Inverse-$n$ Relationship}

The coupling ratio hierarchy---O-001 ($n = 4$, ${\sim}474$K:1) $>$
H-ch6 ($n = 6$, ${\sim}70$K:1) $>$ T-001r2 ($n = 8$,
${\sim}42$K:1)---follows a clear inverse relationship with channel
count (Figure~\ref{fig:inverse_n}). The mechanism is straightforward:
at higher~$n$, the bridge
has more cross-axial entries to suppress, diluting the contrastive
signal per entry. At $n = 4$, only 4~cross-axial pairs compete
with 2~co-axial pairs; at $n = 8$, 24~cross-axial pairs compete
with 4~co-axial pairs. The co-axial/cross-axial pair count ratio
($|\mathcal{C}|/|\mathcal{X}|$) is $2/4 = 0.50$ for octahedral,
$3/12 = 0.25$ for RD, and $4/24 = 0.167$ for tesseract.

This suggests a practical guideline: \emph{use the minimum channel
count that admits the desired topology}. Higher channel counts
provide more topological expressiveness (more independent axes) but
weaker per-axis signal. The practitioner's choice of~$n$ trades off
topological resolution against coupling strength.

\begin{figure}[t]
  \centering
  \includegraphics[width=0.7\textwidth]{figures/fig_inverse_n.pdf}
  \caption{Inverse-$n$ relationship between channel count and
    co-axial/cross-axial coupling ratio. Octahedral ($n = 4$):
    473{,}622:1. RD ($n = 6$): 70{,}404:1. Tesseract ($n = 8$):
    41{,}564:1. Fewer channels concentrate the co-axial signal into
    fewer entries, producing stronger per-pair coupling.}
  \label{fig:inverse_n}
\end{figure}

\subsection{Implications for Adapter Design}

If the bridge is a programmable topology substrate, adapter design
acquires a new degree of freedom. Current practice treats the rank
dimension as homogeneous. Multi-channel extensions introduce
structure but leave it to the optimizer to discover whatever
coupling pattern the task prefers. The Steersman enables
\emph{prescribed} coupling: the practitioner specifies which
channels should interact, and the feedback mechanism enforces the
specification during training.

\begin{itemize}
  \item \textbf{Topology as hyperparameter.} The pair specification
    is now a hyperparameter alongside rank, learning rate, and target
    modules. Different polytope geometries may suit different task
    families.
  \item \textbf{Compositional adapters.} Adapters with known
    block-diagonal structure compose more predictably than dense
    adapters: blocks can be independently swapped, scaled, or merged.
    The $2 \times 2$ block is the natural unit of composition.
  \item \textbf{Multi-modal bridges.} When adapting models across
    modalities (text, image, audio), the pair specification could
    encode known cross-modal relationships in the bridge structure.
  \item \textbf{Diagnostic topology.} Train with $n = 6$ or $n = 8$
    and cybernetic feedback to discover the bridge's preferred
    topology, then deploy with lower~$n$ for parameter efficiency.
    Paper~3 showed that $n = 3$ matches $n = 6$ task performance
    with $4\times$ fewer bridge parameters.
\end{itemize}

\subsection{Zero-Cost Bridge: The P0 Proof}

A practical question precedes all topology programming results: does
the learnable bridge degrade task performance? If the $n^2$ bridge
parameters and Steersman feedback impose a cost on the underlying
fine-tuning objective, topology programming would face a performance
tax before its benefits could be assessed.

We test this directly. On Qwen~2.5-1.5B-Instruct---a different model
family and scale from the TinyLlama experiments above---we train two
configurations with identical hyperparameters (rank~24, seed~42,
alpaca-cleaned, cosine LR schedule): standard LoRA and TeLoRA with a
learnable bridge ($n = 6$, identity initialization, no contrastive
loss).

\begin{table}[H]
\centering
\caption{P0 head-to-head: standard LoRA vs.\ TeLoRA with learnable
  bridge (no contrastive loss). Qwen~2.5-1.5B-Instruct, $r = 24$,
  seed~42, 3{,}235~steps.}
\label{tab:p0}
\small
\begin{tabular}{lcccc}
\toprule
Configuration & Final Loss & Trainable Params & Bridge Params & Time \\
\midrule
Standard LoRA ($r = 24$) & 0.1763 & 3{,}268{,}608 & --- & 42.8 min \\
TeLoRA ($r = 24$, $n = 6$) & 0.1762 & 3{,}270{,}624 & 2{,}016 (0.06\%) & 58.3 min \\
\bottomrule
\end{tabular}
\end{table}

Loss is identical to four significant figures (0.01\% difference,
within noise). The bridge adds 2{,}016~parameters---36~per bridge
matrix across 56~target projections---representing 0.06\% of the
trainable parameter budget. The 36.2\% training time overhead
reflects the Steersman's per-step Fiedler eigenvalue computation
and control law evaluation.

Without contrastive loss, the bridges settle near the spectral
attractor: mean Fiedler~$= 0.038$ (somewhat below the
TinyLlama attractor at~$0.09$, consistent with model-dependent
variation), co-axial/cross-axial ratio~$\approx 1$:1 (isotropic,
no directional preference), and mean deviation from identity~$= 0.07$
(Figure~\ref{fig:p0_proof}). The bridge evolves under spectral
pressure alone but maintains the unprogrammed baseline---exactly
the attractor behavior documented in Section~\ref{sec:attractor}.

\begin{figure}[t]
  \centering
  \includegraphics[width=\textwidth]{figures/fig_p0_proof.pdf}
  \caption{P0 head-to-head comparison on Qwen~2.5-1.5B-Instruct.
    Top left: loss curves overlap exactly.
    Top right: training time comparison (36.2\% overhead from
    Steersman computation). Middle: bridge matrix evolution at
    steps~500, 2{,}000, and~3{,}235---identity-like structure with
    slight evolution but no block-diagonal emergence (no contrastive
    loss applied). Bottom: Fiedler connectivity, co/cross ratio, and
    deviation from identity over training.}
  \label{fig:p0_proof}
\end{figure}

The P0 result establishes that the learnable bridge is a zero-cost
container. It adds negligible parameters, preserves task performance,
and defaults to the spectral attractor when no topology is specified.
Topology programming via the Steersman adds structure to this
container at no additional parameter cost beyond the 36~entries per
bridge.

\subsection{Limitations}

Several limitations qualify the present results.

\textbf{Single task family.} All experiments use instruction-following
fine-tuning on alpaca-cleaned. Whether the four-regime taxonomy,
spectral attractor, and topology programming transfer to other task
types (classification, summarization, code generation) remains
untested.

\textbf{Pair-based specifications only.} The current pair
specification interface partitions channels into pairs. More complex
topologies---cliques of size~${>}\,2$, hierarchical nesting, cyclic
structures---would require extending the contrastive loss
formulation.

\textbf{No formal optimality proof.} We demonstrate that the Steersman
programs a given topology, but we do not prove that any particular
topology is optimal for any downstream task. The topology is a
hyperparameter; its relationship to task performance remains to be
characterized.

\textbf{No theoretical account of the attractor.} The spectral
attractor ($\text{Fiedler} \approx 0.09$) is empirically consistent
across channel counts, but we provide no theoretical derivation
of the equilibrium value. A theoretical account, possibly through
random matrix theory or mean-field analysis of the spectral loss
dynamics, is future work.

\textbf{No theoretical account of geometric selectivity.} We observe
that only polytope-derived pair specifications produce block-diagonal
structure, but we do not provide a formal characterization of which
pair specifications constitute ``valid geometric programs.''
Understanding why the contrastive loss landscape has minima at
polytope-aligned configurations and not at algebraically structured
but spatially unembedded configurations is an open question.

\textbf{Single seed for most experiments.} All experiments use random
seed~42 except Seed-43 and Seed-44 (which use seeds~43 and~44
respectively). While the binary regime classification is robust given
the magnitude of the effects, specific numerical quantities---coupling
ratios, Fiedler values, deviation---are point estimates.

\textbf{Emanation architecture incompletely explored.} E-001 tests one
hierarchical design (master bridge + $1/\sqrt{l}$-scaled offsets). The
space of hierarchical bridge architectures is large; other designs may
produce different outcomes. The third regime (hierarchical coherence)
is characterized by a single experiment.

\textbf{Scale.} All topology experiments use TinyLlama-1.1B. Paper~3
established scale consistency for RD contrastive at 1.1B and 7B;
the generalization of 4D programming, geometric selectivity, and the
four-regime taxonomy to larger models is untested.
